% Original PDF page 1 (printed p. 587).
% Metadata commands are invoked by the main document.
\newcommand{\paperauthors}{%
  A.O. Caldeira\thanks{A.O. Caldeira が1980年9月に University of Sussex に D.Phil. 学位取得のため提出した学位論文の一部に基づく。}\\[0.3em]
  {\small Instituto de Física ``Gleb Wataghin'', Universidade Estadual de Campinas,}\\
  {\small 13.100 Campinas -- SP -- Brasil}\\[0.8em]
  A.J. Leggett\\[0.3em]
  {\small School of Mathematical and Physical Sciences, University of Sussex,}\\
  {\small Brighton BN1 9QH, Sussex, UK}%
}
\newcommand{\originalreceiptdates}{1982年8月25日受理\quad 1983年1月4日改訂}
\newcommand{\paperabstract}{%
Feynman と Vernon の影響汎関数法を、任意の温度における Brownian 運動の研究に適用する。対象とする系と環境との散逸的相互作用について特定の模型を選ぶことにより、影響汎関数を閉じた形で評価し、現象論的な粘性係数など少数のパラメータで表すことができる。$\hbar\to0$ の極限では、影響汎関数の形式から得られる結果が古典的な Fokker--Planck 方程式に帰着することを示す。任意の強さの減衰を受ける単純調和振動子の場合には、任意の温度で、系がある特定の種類の純粋状態から出発するときの完全な密度行列 $\rho(x,x',t)$ の時間発展を明示的に表す式を得る。得られた結果を、量子力学における散逸の問題に対する他の方法の結果と比較する。%
}
